\section{Бор и алгоритм Ахо--Корасик}

На прошлой лекции мы научились искать все вхождения одного шаблона $s$ в текст $t$ за линейное время $O(|s| + |t|)$ с помощью префикс-функции или Z-функции. Теперь обобщим задачу: пусть нам дан целый словарь из $n$ шаблонов $s_1, s_2, \dots, s_n$ и большой текст $t$. Требуется найти вхождения всех строк словаря в текст $t$.

Если запустить алгоритм Кнута--Морриса--Пратта отдельно для каждого шаблона $s_i$, мы потратим
\[
    O(|s_1| + |t|) + O(|s_2| + |t|) + \dots + O(|s_n| + |t|) = O\!\left(\sum_{i=1}^n |s_i| + n \cdot |t|\right)
\]
времени. Когда и текст $t$, и количество шаблонов $n$ велики, слагаемое $n \cdot |t|$ становится неприемлемым: мы $n$ раз перечитываем один и тот же текст. Хочется обработать словарь заранее, а затем пройтись по тексту $t$ всего один раз, уложившись в $O\!\left(\sum |s_i| + |t|\right)$ (плюс размер ответа, если нужно явно выписать все вхождения).

\subsection{Структура данных <<Бор>> (Trie)}

Прежде чем искать строки в тексте, поймём, как компактно хранить словарь и быстро проверять принадлежность строки множеству.

\begin{definition}
    \textbf{Бором} (или \textbf{префиксным деревом}, англ. \textit{trie}) над алфавитом $\Sigma$ называется корневое дерево, удовлетворяющее следующим свойствам:
    \begin{enumerate}
        \item На каждом ребре дерева написан ровно один символ из алфавита $\Sigma$.
        \item Из одной вершины не может выходить двух рёбер с одинаковыми символами.
        \item Некоторые вершины дерева помечены как \textbf{терминальные} ($\text{term}[v] = \text{true}$).
    \end{enumerate}
\end{definition}

Каждой вершине $v$ бора взаимно-однозначно соответствует строка, которая получается конкатенацией символов на пути от корня до $v$. Корень бора соответствует пустой строке $\epsilon$. В дальнейшем мы часто будем отождествлять вершину $v$ со строкой, написанной на пути от корня до неё, и обозначать через $|v|$ глубину вершины (длину этой строки).

Множество строк, хранящихся в боре, "--- это в точности те строки, пути для которых заканчиваются в терминальных вершинах.

\begin{center}
\begin{tikzpicture}[
    scale=1.0,
    vertex/.style={circle, draw=black!80, thick, minimum size=7mm, inner sep=1pt, font=\small},
    term/.style={vertex, fill=NavyBlue!20, double, double distance=1.2pt},
    treeedge/.style={draw=black!80, -Stealth, thick}
]
    \node[vertex] (root) at (0, 3.2) {$0$};
    \node[vertex] (a)    at (-1.8, 2.0) {$1$};
    \node[vertex] (b)    at (1.8, 2.0) {$2$};
    \node[term]   (ab)   at (-1.8, 0.7) {$3$};
    \node[term]   (ba)   at (0.9, 0.7) {$4$};
    \node[vertex] (bb)   at (2.7, 0.7) {$5$};
    \node[vertex] (aba)  at (-1.8, -0.6) {$6$};
    \node[term]   (bba)  at (2.7, -0.6) {$7$};
    \node[term]   (abab) at (-1.8, -1.9) {$8$};

    \draw[treeedge] (root) -- node[above left, inner sep=2pt, font=\small\bfseries, text=NavyBlue] {$\mathtt{a}$} (a);
    \draw[treeedge] (root) -- node[above right, inner sep=2pt, font=\small\bfseries, text=NavyBlue] {$\mathtt{b}$} (b);
    \draw[treeedge] (a) -- node[left, font=\small\bfseries, text=NavyBlue] {$\mathtt{b}$} (ab);
    \draw[treeedge] (ab) -- node[left, font=\small\bfseries, text=NavyBlue] {$\mathtt{a}$} (aba);
    \draw[treeedge] (aba) -- node[left, font=\small\bfseries, text=NavyBlue] {$\mathtt{b}$} (abab);
    \draw[treeedge] (b) -- node[above left, inner sep=2pt, font=\small\bfseries, text=NavyBlue] {$\mathtt{a}$} (ba);
    \draw[treeedge] (b) -- node[above right, inner sep=2pt, font=\small\bfseries, text=NavyBlue] {$\mathtt{b}$} (bb);
    \draw[treeedge] (bb) -- node[right, font=\small\bfseries, text=NavyBlue] {$\mathtt{a}$} (bba);
\end{tikzpicture}
\par\nopagebreak\smallskip
\small\textit{Рис. 2.1. Бор для словаря $\{\mathtt{ab}, \mathtt{abab}, \mathtt{ba}, \mathtt{bba}\}$. Двойным контуром выделены терминальные вершины.}
\end{center}

\subsubsection{Базовые операции со словарём}

Бор поддерживает три стандартные операции контейнера \texttt{std::set<string>}:
\begin{itemize}
    \item \textbf{Поиск строки} $\mathtt{find}(s)$: встаём в корень и поочерёдно считываем символы $s[0]$, $s[1]$, $\dots$, $s[|s|-1]$. На каждом шаге переходим по ребру с текущим символом. Если в какой-то момент нужного ребра нет, строки $s$ в боре точно нет. Если же мы успешно прочитали всю строку и оказались в вершине $v$, остаётся проверить флаг $\text{term}[v]$.
    \item \textbf{Добавление строки} $\mathtt{insert}(s)$: аналогично идём от корня по символам строки $s$. Если ребра по очередному символу $c$ ещё нет, создаём новую вершину и проводим в неё ребро с меткой $c$. В конце пути ставим флаг $\text{term}[v] = \text{true}$.
    \item \textbf{Удаление строки} $\mathtt{erase}(s)$: в простейшем (<<ленивом>>) варианте достаточно дойти до вершины $v$, отвечающей строке $s$, и сбросить флаг $\text{term}[v] = \text{false}$. Если же мы хотим реально освобождать память от <<мёртвых>> веток, то после сброса флага нужно подниматься от $v$ к корню по ссылкам на предков и удалять вершины, пока текущая вершина нетерминальная и из неё не выходит ни одного ребра.
\end{itemize}

\begin{note}
    Если в словаре могут встречаться одинаковые строки (мультимножество), вместо булева флага $\text{term}[v]$ в вершине хранят целое число "--- сколько раз данная строка была добавлена в бор, либо список номеров строк словаря, заканчивающихся в $v$.
\end{note}

\subsubsection{Способы хранения переходов и оценки памяти}

Пусть в бор добавлены строки $s_1, \dots, s_n$. Поскольку при добавлении строки $s_i$ мы создаём не более $|s_i|$ новых вершин, общее число вершин $|V|$ удовлетворяет оценке
\[
    |V| \le 1 + \sum_{i=1}^n |s_i|.
\]
Время работы операций и расход памяти зависят от того, как именно в каждой вершине хранятся переходы по символам алфавита $\Sigma$:

\begin{center}
\begin{tabular}{@{}lccc@{}}
    \toprule
    \textbf{Хранение переходов} & \textbf{Память} & \textbf{Переход} & \textbf{Операция с $s$} \\
    \midrule
    Массив длины $|\Sigma|$ (\texttt{vector<int>}) & $O(|V| \cdot |\Sigma|)$ & $O(1)$ & $O(|s|)$ \\
    Дерево поиска (\texttt{std::map}) & $O(|V|)$ & $O(\log |\Sigma|)$ & $O(|s| \log |\Sigma|)$ \\
    Хеш-таблица (\texttt{unordered\_map}) & $O(|V|)$ & $O(1)$ в среднем & $O(|s|)$ в среднем \\
    \bottomrule
\end{tabular}
\end{center}

На практике при небольшом фиксированном алфавите (например, строчные латинские буквы, $|\Sigma| = 26$) почти всегда используют массив фиксированной длины: он работает быстрее всего и не тратит лишней памяти. Если же алфавит большой, в вершинах хранят \texttt{std::map} (тогда переход из вершины $v$ с $k_v \le |\Sigma|$ исходящими рёбрами работает за $O(\log k_v)$) или хеш-таблицу \texttt{std::unordered\_map}: суммарное число рёбер во всём дереве равно $|V| - 1$, поэтому суммарный размер всех словарей переходов будет $O(|V|)$ независимо от мощности алфавита.

\begin{lstlisting}[caption={Реализация бора с массивом переходов фиксированного размера}]
const int SIGMA = 26;

struct Node {
    vector<int> next;
    bool term = false;

    Node() : next(SIGMA, -1) {}
};

struct Trie {
    vector<Node> nodes;

    Trie() {
        nodes.emplace_back(); // корень имеет индекс 0
    }

    void insert(const string& s) {
        int v = 0;
        for (char ch : s) {
            int c = ch - 'a';
            if (nodes[v].next[c] == -1) {
                nodes[v].next[c] = nodes.size();
                nodes.emplace_back();
            }
            v = nodes[v].next[c];
        }
        nodes[v].term = true;
    }

    bool find(const string& s) const {
        int v = 0;
        for (char ch : s) {
            int c = ch - 'a';
            if (nodes[v].next[c] == -1) {
                return false;
            }
            v = nodes[v].next[c];
        }
        return nodes[v].term;
    }
};
\end{lstlisting}

\begin{note}[Числовой бор]
    Бор отлично подходит не только для строк, но и для целых чисел. Любое $W$-битное целое число (например, типа \texttt{uint32\_t} или \texttt{uint64\_t}) можно рассматривать как бинарную строку длины $W$, записанную от старшего бита к младшему. В таком бинарном боре $|\Sigma| = 2$, а все терминальные вершины лежат строго на глубине $W$.

    Классическое применение числового бора "--- задача о максимальном $\operatorname{XOR}$: дано множество чисел $a_1, \dots, a_n$ и запрос $x$, требуется найти $\max_i (a_i \oplus x)$. Спускаясь от корня бора по битам числа $x$ от старшего к младшему, мы на каждом шаге жадно пытаемся пойти по ребру с битом, противоположным текущему биту $x$. Если такое ребро есть, мы гарантированно получаем единицу в соответствующем разряде ответа; если нет "--- вынужденно идём по совпадающему биту. Время ответа на запрос составляет $O(W)$.
\end{note}

\subsection{Алгоритм Ахо--Корасик}

Вернёмся к задаче поиска набора шаблонов $s_1, \dots, s_n$ в тексте $t$. Мы разберём две её классические постановки:
\begin{enumerate}
    \item \textbf{Подсчёт суммарного числа вхождений}: требуется найти одно число "--- сколько всего раз строки из словаря встречаются в тексте $t$ (с учётом кратности).
    \item \textbf{Перечисление всех вхождений}: для каждой позиции текста $t$ требуется явно вывести все шаблоны $s_i$, которые в ней заканчиваются (именно эту задачу решает стандартная консольная утилита \texttt{grep}).
\end{enumerate}

Построим бор из всех шаблонов $s_1, \dots, s_n$. Что вообще означает, что строка $s_i$ входит в текст $t$ и заканчивается в позиции $j$? Это значит, что $s_i$ является суффиксом префикса $t[0 \dots j]$.

Будем читать текст $t$ посимвольно слева направо. Представим, что мы прочитали префикс текста $\tau = t[0 \dots j]$. Нам хочется знать, какие суффиксы строки $\tau$ являются терминальными вершинами нашего бора. Основная идея алгоритма Ахо--Корасик состоит в следующем: вместо того чтобы отслеживать все суффиксы строки $\tau$ в боре одновременно, будем поддерживать \textbf{только самый длинный суффикс} строки $\tau$, который является вершиной бора.

Почему этой информации достаточно? Пусть $v$ "--- самый длинный суффикс строки $\tau$, присутствующий в боре. Тогда любой другой суффикс строки $\tau$, который тоже есть в боре, обязан быть суффиксом строки $v$ (ведь из двух суффиксов одной и той же строки $\tau$ более короткий всегда является суффиксом более длинного). Значит, зная вершину $v$, мы однозначно знаем и все остальные подходящие суффиксы "--- они зависят только от самой вершины $v$ в боре, но не от предыстории текста $\tau$.

Чтобы быстро пересчитывать текущую вершину $v$ при дописывании очередного символа текста $c = t[j+1]$, превратим наш бор в конечный автомат.

\begin{definition}
    Пусть $v$ "--- некоторая вершина бора.
    \begin{itemize}
        \item \textbf{Суффиксной ссылкой} $\link(v)$ назовём вершину бора, отвечающую самому длинному \emph{собственному} суффиксу строки $v$, который присутствует в боре (то есть $\link(v) \ne v$). Для корня бора $0$ собственный суффикс не существует, поэтому положим $\link(0) = -1$.
        \item \textbf{Автоматным переходом} $\step(v, c)$ по символу $c \in \Sigma$ назовём вершину бора, отвечающую самому длинному суффиксу строки $v + c$, который присутствует в боре.
    \end{itemize}
\end{definition}

Заметим тонкое различие в определениях: в $\link(v)$ мы требуем \emph{собственный} суффикс (поэтому $|\link(v)| < |v|$), а в $\step(v, c)$ ограничений на собственность нет "--- если в боре есть прямое ребро $v \xrightarrow{c} u$, то вся строка $v + c$ лежит в боре и $\step(v, c) = u$.

\begin{center}
\begin{tikzpicture}[
    scale=1.0,
    vertex/.style={circle, draw=black!80, thick, minimum size=7mm, inner sep=1pt, font=\small},
    term/.style={vertex, fill=NavyBlue!20, double, double distance=1.2pt},
    treeedge/.style={draw=black!80, -Stealth, thick},
    linkedge/.style={draw=BrickRed, -Stealth, dashed, thick}
]
    \node[vertex] (0) at (0, 3) {$0$};
    \node[vertex] (1) at (-2, 1.6) {$1$};
    \node[vertex] (2) at (2, 1.6) {$2$};
    \node[term]   (3) at (-2, 0) {$3$};
    \node[term]   (4) at (1, 0) {$4$};
    \node[vertex] (5) at (3, 0) {$5$};
    \node[vertex] (6) at (-2, -1.6) {$6$};
    \node[term]   (7) at (3, -1.6) {$7$};

    \draw[treeedge] (0) -- node[below right, inner sep=2pt, font=\small\bfseries, text=NavyBlue] {$\mathtt{a}$} (1);
    \draw[treeedge] (0) -- node[below left, inner sep=2pt, font=\small\bfseries, text=NavyBlue] {$\mathtt{b}$} (2);
    \draw[treeedge] (1) -- node[left, font=\small\bfseries, text=NavyBlue] {$\mathtt{b}$} (3);
    \draw[treeedge] (2) -- node[left, font=\small\bfseries, text=NavyBlue] {$\mathtt{a}$} (4);
    \draw[treeedge] (2) -- node[below left, inner sep=2pt, font=\small\bfseries, text=NavyBlue] {$\mathtt{b}$} (5);
    \draw[treeedge] (3) -- node[left, font=\small\bfseries, text=NavyBlue] {$\mathtt{a}$} (6);
    \draw[treeedge] (5) -- node[right, font=\small\bfseries, text=NavyBlue] {$\mathtt{a}$} (7);

    % Suffix links
    \draw[linkedge, bend left=22] (1) to (0);
    \draw[linkedge, bend right=22] (2) to (0);
    \draw[linkedge, bend left=18] (3) to (2);
    \draw[linkedge] (4) to (1);
    \draw[linkedge, bend right=25] (5) to (2);
    \draw[linkedge] (6) to (4);
    \draw[linkedge] (7) to (4);
\end{tikzpicture}
\par\nopagebreak\medskip
\small\textit{Рис. 2.2. Суффиксные ссылки (показаны красным пунктиром) в боре для $\{\mathtt{ab}, \mathtt{aba}, \mathtt{ba}, \mathtt{bba}\}$. Например, для вершины $6$ ($\mathtt{aba}$) наибольшим собственным суффиксом в боре является вершина $4$ ($\mathtt{ba}$).}
\end{center}

\subsubsection{Рекуррентные формулы и построение обходом в ширину}

Как быстро вычислить $\link(v)$ и $\step(v, c)$ для всех вершин бора? Оказывается, они взаимно выражаются друг через друга по простым формулам.

\begin{lemma}
    Пусть $\nxt(v, c)$ обозначает сына вершины $v$ по ребру с символом $c$ в боре (или $-1$, если такого ребра нет). Тогда справедливы следующие соотношения:
    \begin{enumerate}
        \item \textbf{База для корня} $v = 0$: $\link(0) = -1$, а для любого символа $c \in \Sigma$
        \[
            \step(0, c) = \begin{cases}
                \nxt(0, c), & \text{если } \nxt(0, c) \ne -1, \\
                0, & \text{иначе.}
            \end{cases}
        \]
        \item \textbf{Пересчёт суффиксной ссылки}: пусть в боре есть ребро $v \xrightarrow{c} u$ (то есть $u = \nxt(v, c)$). Тогда
        \[
            \link(u) = \begin{cases}
                0, & \text{если } v = 0, \\
                \step(\link(v), c), & \text{если } v \ne 0.
            \end{cases}
        \]
        \item \textbf{Пересчёт автоматного перехода}: для любой вершины $u \ne 0$ и символа $d \in \Sigma$
        \[
            \step(u, d) = \begin{cases}
                \nxt(u, d), & \text{если } \nxt(u, d) \ne -1, \\
                \step(\link(u), d), & \text{иначе.}
            \end{cases}
        \]
    \end{enumerate}
\end{lemma}

\begin{proof}
    \textbf{1.} Для корня $0$ строка пуста. При дописывании символа $c$ мы получаем однобуквенную строку $c$. У неё всего два суффикса: сама строка $c$ (длины $1$) и пустая строка $\epsilon$ (длины $0$). Если из корня есть ребро по $c$, то строка $c$ присутствует в боре и является ответом; иначе ответом служит корень $0$.

    \textbf{2.} Пусть $v \xrightarrow{c} u$, так что строка $u$ имеет вид $v + c$. Если $v = 0$, то $|u| = 1$, и единственным собственным суффиксом строки $u$ является пустая строка $0$, поэтому $\link(u) = 0$.

    Пусть теперь $v \ne 0$. Нам нужно найти самый длинный собственный суффикс строки $v + c$, лежащий в боре. Рассмотрим строку $w = \link(v)$ "--- самый длинный собственный суффикс строки $v$ в боре. По определению, $\step(\link(v), c)$ есть самый длинный суффикс строки $w + c$ в боре. Поскольку $w$ "--- собственный суффикс $v$, строка $w + c$ является собственным суффисом $v + c = u$.

    Почему в боре не может найтись ещё более длинного собственного суффикса строки $u$? Предположим противное: пусть в боре есть собственный суффикс $\alpha$ строки $u$ такой, что $|\alpha| > |\step(\link(v), c)| \ge 0$. Так как $\alpha$ непуст и является суффиксом $v + c$, он заканчивается на букву $c$, то есть $\alpha = \alpha' + c$, где $\alpha'$ "--- собственный суффикс строки $v$. Поскольку вершина $\alpha$ есть в боре, её префикс $\alpha'$ также обязан быть вершиной бора. Но если $|\alpha'| \le |\link(v)|$, то $\alpha'$ является суффиксом $\link(v)$, а значит, $\alpha = \alpha' + c$ является суффиксом $\link(v) + c$ в боре, и его длина не может превосходить $|\step(\link(v), c)|$. Следовательно, $|\alpha'| > |\link(v)|$, что противоречит определению $\link(v)$ как \emph{самого длинного} собственного суффикса $v$ в боре.

    \begin{center}
    \begin{tikzpicture}[
        scale=1.0,
        vertex/.style={circle, draw=black!80, thick, minimum size=6.5mm, inner sep=1pt, font=\small},
        treeedge/.style={draw=black!80, -Stealth, thick},
        linkedge/.style={draw=BrickRed, -Stealth, dashed, thick},
        stepedge/.style={draw=OliveGreen!80!black, -Stealth, dashdotted, thick}
    ]
        \node[vertex] (root) at (0, 2.2) {$0$};
        \node[vertex] (v) at (-2.4, 0.7) {$v$};
        \node[vertex] (u) at (-2.4, -1.0) {$u$};
        \node[vertex] (lv) at (1.6, 0.9) {$w$};
        \node[vertex] (lu) at (2.4, -1.0) {$z$};

        \draw[treeedge, decorate, decoration={snake, amplitude=1pt, segment length=8pt, post length=3pt}] (root) -- (v);
        \draw[treeedge, decorate, decoration={snake, amplitude=1pt, segment length=8pt, post length=3pt}] (root) -- (lv);
        \draw[treeedge] (v) -- node[left, font=\small\bfseries, text=NavyBlue] {$c$} (u);
        \draw[linkedge] (v) -- node[above, font=\footnotesize, text=BrickRed] {$\link(v)$} (lv);
        \draw[stepedge] (lv) -- node[right, font=\footnotesize, text=OliveGreen!80!black] {$\step(w, c)$} (lu);
        \draw[linkedge] (u) -- node[below=3pt, font=\footnotesize, text=BrickRed] {$\link(u) = \step(\link(v), c)$} (lu);
    \end{tikzpicture}
    \end{center}

    \textbf{3.} Рассуждение для $\step(u, d)$ совершенно аналогично: если из $u$ в боре есть ребро по $d$, то вся строка $u + d$ присутствует в боре и равна $\nxt(u, d)$. Если же ребра нет, то искомый наибольший суффикс строки $u + d$ является собственным, а значит, имеет вид $\beta + d$, где $\beta$ "--- собственный суффикс $u$ в боре. Самый длинный такой $\beta$ есть $\link(u)$, поэтому ответ совпадает с $\step(\link(u), d)$.
\end{proof}

Обратим внимание на глубины вершин в доказанных формулах: так как $\link(v)$ "--- собственный суффикс $v$, выполнено строгое неравенство $|\link(v)| < |v|$. Значит:
\begin{itemize}
    \item для расчёта $\link(u)$ нужны $\link(v)$ и переход $\step(\link(v), c)$, где $|\link(v)| < |v| < |u|$;
    \item для расчёта $\step(u, d)$ нужны $\link(u)$ и переход $\step(\link(u), d)$, где $|\link(u)| < |u|$.
\end{itemize}
Иными словами, значения $\link$ и $\step$ для любой вершины зависят только от вершин \textbf{строго меньшей глубины}. Поэтому мы можем вычислить их за один проход, запустив \textbf{обход в ширину (BFS)} из корня бора: ведь BFS посещает вершины дерева в порядке неубывания их глубины.

\begin{lstlisting}[caption={Построение автомата Ахо--Корасик обходом в ширину}]
struct Node {
    vector<int> next;
    vector<int> step;
    int link = -1;
    bool term = false;
    int cnt = 0; // число словарных строк, являющихся суффиксами данной вершины

    Node() : next(SIGMA, -1), step(SIGMA, -1) {}
};

void init(vector<Node>& nodes) {
    nodes[0].link = -1;
    for (int c = 0; c < SIGMA; ++c) {
        if (nodes[0].next[c] != -1) {
            nodes[0].step[c] = nodes[0].next[c];
        } else {
            nodes[0].step[c] = 0;
        }
    }

    queue<int> q;
    q.push(0);

    while (!q.empty()) {
        int v = q.front();
        q.pop();

        // Проталкиваем информацию сыновьям вершины v
        for (int c = 0; c < SIGMA; ++c) {
            int u = nodes[v].next[c];
            if (u == -1) continue;

            // 1. Вычисляем суффиксную ссылку для сына u
            if (v == 0) {
                nodes[u].link = 0;
            } else {
                nodes[u].link = nodes[nodes[v].link].step[c];
            }

            // 2. Вычисляем автоматные переходы из сына u по всем буквам d
            for (int d = 0; d < SIGMA; ++d) {
                if (nodes[u].next[d] != -1) {
                    nodes[u].step[d] = nodes[u].next[d];
                } else {
                    nodes[u].step[d] = nodes[nodes[u].link].step[d];
                }
            }

            // 3. Предподсчёт числа терминальных вершин на пути суффиксных ссылок
            nodes[u].cnt = nodes[nodes[u].link].cnt + (nodes[u].term ? 1 : 0);

            q.push(u);
        }
    }
}
\end{lstlisting}

\begin{theorem}
    Построение автомата Ахо--Корасик по бору с $|V|$ вершинами над алфавитом $\Sigma$ работает за время $O(|V| \cdot |\Sigma|)$ и требует $O(|V| \cdot |\Sigma|)$ памяти.
\end{theorem}

\begin{proof}
    Каждая вершина $v$ попадает в очередь BFS ровно один раз. Во внутреннем цикле мы перебираем только существующие в дереве рёбра $v \xrightarrow{c} u$ "--- суммарно по всем вершинам таких рёбер ровно $|V| - 1$. Для каждого ребра $(v, u)$ мы за $O(1)$ вычисляем $\link(u)$ и за $O(|\Sigma|)$ заполняем массив $\step(u, d)$ для всех $d \in \Sigma$. Следовательно, суммарное время работы составляет $O(|V| \cdot |\Sigma|)$.
\end{proof}

\subsection{Решение задачи о вхождениях и сжатые суффиксные ссылки}

Теперь, когда автомат построен, посмотрим, как за один проход по тексту $t$ решить обе постановки нашей задачи.

\subsubsection{Постановка 1: суммарное количество вхождений}

Пусть после прочтения префикса $t[0 \dots j]$ мы находимся в вершине автомата $v$. Какие словарные строки заканчиваются в позиции $j$? Любая такая строка является суффиксом $t[0 \dots j]$ и лежит в боре, а значит, является суффиксом вершины $v$.

Самый длинный собственный суффикс $v$ в боре "--- это $\link(v)$; следующий по длине "--- $\link(\link(v))$, и так далее вплоть до корня $0$. Никаких других суффиксов строки $v$ в боре нет. Значит, все словарные строки, заканчивающиеся в текущей позиции текста, взаимно-однозначно соответствуют терминальным вершинам на цепочке суффиксных ссылок
\[
    v \to \link(v) \to \link(\link(v)) \to \dots \to 0.
\]
Обозначим через $\text{cnt}[v]$ количество терминальных вершин на этом пути (если в словаре были одинаковые шаблоны, в каждой вершине суммируем кратность её терминальности). Так как путь из $u$ получается добавлением вершины $u$ к пути из $\link(u)$, величина $\text{cnt}[u]$ пересчитывается по формуле
\[
    \text{cnt}[u] = \text{cnt}[\link(u)] + [\text{term}[u] = \text{true}]
\]
прямо внутри BFS при построении автомата (см. пункт 3 во внутреннем цикле листинга выше).

После этого поиск суммарного числа вхождений в текст $t$ сводится к простому циклу:
\begin{lstlisting}[caption={Подсчёт суммарного числа вхождений за $O(|t|)$}]
long long total_occurrences(const vector<Node>& nodes, const string& t) {
    int v = 0;
    long long ans = 0;
    for (char ch : t) {
        v = nodes[v].step[ch - 'a'];
        ans += nodes[v].cnt;
    }
    return ans;
}
\end{lstlisting}
Итоговая асимптотика решения первой задачи равна $O(|V| \cdot |\Sigma| + |t|)$.

\subsubsection{Постановка 2: вывод всех вхождений и сжатые ссылки}

Что изменится, если нам нужно не просто посчитать количество вхождений, а явно перечислить все шаблоны, заканчивающиеся в каждой позиции $j$ текста $t$?

Если мы в каждой позиции $j$ будем честно подниматься по цепочке суффиксных ссылок $v, \link(v), \link(\link(v)), \dots$ до корня и проверять флаг $\text{term}$, то время работы может оказаться квадратичным. Действительно, на пути от $v$ до корня может встретиться длинная цепочка из сотен \emph{нетерминальных} вершин, по которым мы будем впустую шагать на каждом символе текста, не находя ни одного вхождения.

Чтобы пропускать нетерминальные вершины за $O(1)$, сожмём эти пути.

\begin{definition}
    \textbf{Сжатой суффиксной ссылкой} (или \textbf{терминальной ссылкой}) $\compressed(v)$ вершины $v$ назовём ближайшую к $v$ терминальную вершину на пути по суффиксным ссылкам $\link(v), \link(\link(v)), \dots$ (не считая самой вершины $v$). Если на этом пути нет ни одной терминальной вершины, положим $\compressed(v) = -1$.
\end{definition}

Сжатая ссылка легко вычисляется во время того же BFS по правилу:
\[
    \compressed(u) = \begin{cases}
        -1, & \text{если } u = 0, \\
        \link(u), & \text{если } \text{term}[\link(u)] = \text{true}, \\
        \compressed(\link(u)), & \text{иначе.}
    \end{cases}
\]
Теперь, находясь в состоянии $v$ после чтения $t[j]$, мы сначала проверяем саму вершину $v$ (если $\text{term}[v] = \text{true}$, выводим номер её строки), а затем прыгаем по сжатым ссылкам $v \la \compressed(v)$, пока не дойдём до $-1$. Поскольку каждый прыжок по сжатой ссылке гарантированно приводит нас в терминальную вершину и даёт хотя бы одно новое вхождение в ответ, количество прыжков в точности равно размеру ответа $\text{ans}$.

Следовательно, вторая постановка задачи решается за оптимальное время $O(|V| \cdot |\Sigma| + |t| + \text{ans})$.

\subsection{Применение: декодирование кода без разделителей}

\begin{problem}[Декодирование кода без пауз]
    Каждой букве $i$-го типа ($1 \le i \le n$) сопоставлен свой непустой бинарный код $s_i \in \{0, 1\}^+$. Все коды $s_i$ попарно различны, сумма их длин равна $L = \sum_{i=1}^n |s_i|$. Дана строка $t \in \{0, 1\}^*$, полученная записью кодов некоторой последовательности букв подряд без пробелов. Требуется определить, можно ли декодировать строку $t$ (то есть разбить её на словарные коды), и если да "--- восстановить хотя бы одно разбиение.
\end{problem}

\begin{solution}
    Решим задачу методом динамического программирования. Пусть $\text{dp}[j]$ ($0 \le j \le |t|$) "--- булева величина, равная $\text{true}$, если префикс $t[0 \dots j-1]$ длины $j$ можно разбить на кодовые слова, и $\text{false}$ иначе. База динамики: $\text{dp}[0] = \text{true}$.

    Для вычисления $\text{dp}[j]$ при $j \ge 1$ нам нужно перебрать последнее кодовое слово в разбиении префикса $t[0 \dots j-1]$. Это кодовое слово $s_i$ должно заканчиваться в позиции $j-1$ (то есть быть суффиксом $t[0 \dots j-1]$), а оставшийся префикс длины $j - |s_i|$ должен быть корректно декодируемым ($\text{dp}[j - |s_i|] = \text{true}$).

    Построим автомат Ахо--Корасик по словарю кодов $\{s_1, \dots, s_n\}$ со сжатыми суффиксными ссылками. Будем идти по строке $t$ слева направо, поддерживая текущее состояние автомата $v$. В позиции $j$ все кодовые слова $s_i$, являющиеся суффиксами $t[0 \dots j-1]$, мы находим, прыгая из $v$ по сжатым ссылкам $\compressed$. Для каждого найденного терминального состояния с длиной кода $\ell = |s_i|$, если $\text{dp}[j - \ell] = \text{true}$, мы полагаем $\text{dp}[j] = \text{true}$ и запоминаем индекс $i$ для восстановления ответа.

    \begin{exercise}
        Покажите, что в любой позиции $j$ число различных словарных строк, заканчивающихся в $j$, не превосходит $\min(n, \sqrt{2L})$. \textit{(Указание: терминальные вершины на пути суффиксных ссылок имеют попарно различные длины $\ell_1 < \ell_2 < \dots < \ell_k$, причём $\ell_1 + \ell_2 + \dots + \ell_k \le L$).} Отсюда следует, что время работы описанного алгоритма составляет $O(L + |t| \sqrt{L})$.
    \end{exercise}
\end{solution}
